% !TeX root = ../papers/nc-p3.tex
\section{空間有理正規三次の六点残余と四移動}
\label{sec:twisted-cubic-residual}
四辺で左右の点射影ファイバーが空か一点となる標準有限Frobenius窓で、
点グラフが空間有理正規三次曲線 $C$ と交わらない零次元部分 $F$ の和となる場合を扱う。
各文字空間は同じ $\mathbb P^1$ の完全な二項三次形式の空間として同一視できる。
\begin{theorem}\label{thm:twisted-cubic-residual}
辺の六関係を二項変数対で表し、対角の一次交代因子で割ると、
二次形式二組の九次元積空間内の六平面 $T$ を得る。
$\Pi=\mathbb P(T^\perp)$ の行列座標を $B=(B_{ij})_{0\le i,j\le2}$ とする。
$F$ はこの平面での次の二式の完全交叉とスキームとして同型である。
\[
 \det B=0,\qquad
 Q=B_{00}B_{22}-(B_{01}+B_{10})(B_{12}+B_{21})
       + B_{11}(B_{02}+B_{20}+B_{11})=0 .
\]
この長さ6の交叉の全ての点で $B$ は階数2である。
\end{theorem}
\begin{proof}
曲線対の基点を一致させた座標で、六関係を対角上で零となる二次数 $(3,3)$ 形式とする。
対角因子による除算後の全九形式の双対写像は、二点 $X,Y\in\mathbb P^3$ に対し
\[
 B_{ij}=X_iY_{j+1}-X_{i+1}Y_j,\qquad 0\le i,j\le2
\]
を与える。その基点は曲線の対角である。実際、二つの三行二列Hankel行列
$H(X)=(X_i,X_{i+1})$ と $H(Y)$ の間に非退化交代行列を挟んだ積が $B$ である。
両方が階数2なら $B$ も階数2で、一方だけが階数2でも零にならない。
両方階数1なら二点は曲線上で、同じ点でだけ零となる。
基点は被約でもある。曲線の点を $X=Y=[1:0:0:0]$ とすると、
一次部は $Y_1-X_1,Y_2,Y_3,-X_2,-X_3$ を含み余次元5の法方向を全て切る。
二項変数の同時基底変更で全点に運べる。

写像の像の閉包は二式の完全交叉である。行列積による行列式消滅と、
$X\wedge Y$ のPlücker関係による $Q$ の消滅は表示に代入しても確かめられる。
対称成分の六座標は $X\wedge Y$ の六座標を可逆線型変換したものである
（$B_{ii}$ は隣接小行列、$B_{ij}+B_{ji}$ で残りを順に得る）。
階数2行列の左右核を二項二次式として射影平面二つに写す。
各核式の零因子は長さ2であり、その因子の張る三次曲線の割線上を $X,Y$ が動くと、
同じ両核を持つ行列の射影三空間内で写像が得られる。
一般の両因子の台 $p_1,p_2$ と $q_1,q_2$ が全て異なれば、割線の座標で
写像はSegre写像の四成分にそれぞれ $[p_i,q_j]$ を掛けるだけである（角括弧は二次元の行列式）。
従ってこのファイバーでは既約な二次曲面をちょうど一回被覆する。

余分な六次元成分が二核の特殊部分に映ることもない。階数1以下の行列は射影次元4である。
階数2では核の各ファイバーの次元は3以下なので、基が余次元2以上なら足りない。
基の因子上で $Q$ が恒等零になるものはない。実際、相異なる単純根で台が交わらなければ既述の二次曲面である。
片方に二重根があっても両因子が交わらないなら、長さ2の各座標環のテンソル積上で
角括弧は単元であり、同じSegre座標の可逆変換となる。ここで $Q$ 自身がその単元倍の式であることは、
二つの割線が交わらないため、$X\wedge Y$ による四交差座標の線型写像が単射で、
Plücker式のその部分空間への制限が通常の二行列式になることからも従う。
単純根が一つだけ一致する場合には、両核で定めた四座標での $Q$ は非零の二座標の積である。
具体的には評価ベクトル $v_2(p_i)v_2(q_j)^t$ による座標を $b_{ij}$ とすると、
$Q$ は $[p_1,q_2][p_2,q_1]b_{11}b_{22}-[p_1,q_1][p_2,q_2]b_{12}b_{21}$ の非零倍である。
比例係数は両対の内部の行列式だけからなる単元である（$Q$ の式に四つの評価ベクトルを代入する）。
従って二式の零点は既約、一般点で被約で、完全交叉の純性からスキームとしてこの像閉包になる。

基点の曲線を吹き上げて写像を延長する。例外因子上の像は、基点での九式の法微分の像を射影化したものである。
$\Pi$ はその像に交わらない。交われば六関係の微分が一つの非零法方向で全て零になり、
滑らかで孤立部と交わらない曲線の接空間が二次元以上になってしまう。
ゆえに $\Pi$ 上の交叉は全て実際の曲線外の点対から来る。その逆像は $F$ なので交叉は有限、次数6である。
これらの点対はどちらの座標も曲線外であるため $B$ は階数2である。
二核に対応する二割線が同じなら、その割線の積上で $B$ の像は高々直線となる。
実際、同じ二根なら角括弧による乗法の階数は2以下である
（二重根では双対数二組の一様化元の差の乗法が階数2）。よって曲線外の点対を通るファイバーが正次元となり矛盾する。
二割線は異なり、$X\ne Y$ である。$B$ が決める $X\wedge Y$ の直線を各割線と交わらせれば $X,Y$ を復元できる。
この直線が割線の一方に等しければ、他方の点が両割線の交点になる。
有理正規三次で異なる二割線が交わるのは因子の共通台の曲線点だけである（長さ4まで線型独立）。
よって交点構成は各点の近くで定数階数の線型代数により正則である。
被約な像閉包のこの開部分でこれはもとの写像の逆であり、平面切断にも同型が残る。
\end{proof}

\begin{theorem}\label{thm:twisted-cubic-monodromy}
この三次曲線上の標準四移動 $\sigma$ は恒等または対合である。
窓テンソルは四つの隣接対角因子を掛けた
二次元文字四つのテンソル $f$（係数16）で与えられる。
対合の場合は二固定点を基底に取ると $f$ は偶数個または奇数個の第二座標を持つ八成分だけに支持される。
\end{theorem}
\begin{proof}
各点写像を $j_i$ とし、$j_{i+4}=j_i\circ\sigma$ の規約で曲線座標をそろえる。
四次テンソル $w$ を四組の二項変数で各次数3の形式として見る。
任意の隣接対をそのグラフの点で評価すると残り二位置の全係数が零になる。
従って $w$ は四つの既約な対角因子の積で割り切れ、商 $f$ は各次数1である。
最後の対角にだけ $\sigma$ または座標同一視の逆向きの $\sigma^{-1}$ が入る。
標準四移動は四次交叉直線を保つので $f$ は四位置同時の $\sigma$ の固有ベクトルとなる。

非自明なユニポテントなら、$\sigma$ のべきの閉包は連結加法群である。
長さ6の非空な有限部の台の各点はこの群で固定されねばならない。
しかし有理正規三次の文字表現にはこの群の射影固定点は曲線上の一点しかない。

半単純の場合は各点の座標 $X_0,\ldots,X_3$ の重みを共通倍率を除いて $0,1,2,3$ とする。
位数が5以上または無限なら、点空間の曲線外の固定点は第二・第三座標点だけとなる。
また $B_{ij}$ の重みは共通倍率を除いて $i+j$ である。重み0と4の純行列は階数1、
重み1と3の純行列はともに $Q=\det=0$ を満たす（それぞれ単一の斜め列だけに支持）。
従って $\Pi$ に重み $\ne2$ の固有ベクトルがあれば平面切断に固定点が生じる。
前定理の復元により、それは曲線外の固定点の対でなければならず、
かつ二つの二次核は異なる。よって第二・第三座標点を互いに結ぶ二通りだけであり、その行列の重みは2である。
従って $\Pi$ は重み2の平面全体に等しい。しかしこの平面の $B_{11}=0$ の直線は二式で消え、有限性に反する。
この議論は窓が対角トーラス全体を保つ場合にも適用できる。

位数4と3を調べる。アフィン変数を $z_0,\ldots,z_3$ とし、順に隣接因子を
$z_1-z_0,z_2-z_1,z_3-z_2,u z_0-z_3$ とする（$u$ は原始根）。
もし $f$ が整数重み一つに支持されれば窓を対角トーラスが保つため既に除外した。
位数4なら従って $f$ は定数項と四変数の積だけを持つ。
隣接切断 $01|23$ の行列は、その両側の対角因子を除いた
$f(z_2-z_1)(u z_0-z_3)$ の係数行列（各側九次元）で階数が決まる。
位置0,1で指数対 $(0,2),(2,0)$ の行はともに零なので、
$\Pi$ に $B_{02},B_{20}$ だけを動かす直線が入り、再び有限性に反する。

位数3なら二固定点を交換して、定数項1と、三次単項式
$\prod_{j\ne m}z_j$ の係数 $A,B,C,D$ だけを仮定できる。
上の九行九列を行の全指数を法3として分けると、三ブロックは列順序を辞書順として
\[
 \begin{pmatrix}0&-1&0\\ C&D-Au&0\\0&(C-B)u&Du\end{pmatrix},\quad
 \begin{pmatrix}1&0&-A\\0&u&-B\\-Cu&-Du&0\end{pmatrix},\quad
 \begin{pmatrix}0&A&0\\-u&B-C&Au-D\\0&0&Bu\end{pmatrix}.
\]
総階数は $4+\#\{A,B,C,D\text{ の非零係数}\}+\mathbf1_{ACu+BD\ne0}$ である。
標準の階数6になるには非零係数はちょうど二つで、それらは反対位置ではない。
すると、係数が零の隣接二位置では、$f$ の全三次項がその二変数の両方を含む。
この二位置を一方にした切断でも、定数項から各変数の指数は高々1、
三次項からは両指数が少なくとも1となる。従って $(0,2),(2,0)$ の行が零で同じ矛盾となる。
これで残るのは恒等または位数2だけで、対合の八成分の記述はその固有指標から従う。
\end{proof}
